\section{IMÁGENES DE REFERENCIA}
\subsection{DISPONIBILIDAD DE LOS DATOS}
\begin{table}[H]
	\begin{tblr}{
    	row{2} = {bg=SuCielo},
        colspec = {X}
	}
    	DISPONIBILIDAD DE LOS DATOS\\
    	ACELERÓGRAFO EN TERRENO NATURAL\\
    	\SetCell{c}\includegraphics[width=17cm, trim=173 100 137 80, clip=true]{example-image}\\
    	No se encontró perdida en la información en el período descargado [\DiaIni/\MesIni/\AnioIni~al \DiaMantto/\MesMantto/\AnioMantto]\\
	\end{tblr}
\end{table}
%
% \begin{table}[H]
% 	\begin{tblr}{
%     	row{2} = {bg=RoyalBlue!50},
%     	row{4} = {bg=gray!25},
%         colspec = {X}
% 	}
%     	DISPONIBILIDAD DE LOS DATOS\\
%     	ACELERÓGRAFO AZOTEA\\
%     	\SetCell{c}\includegraphics[width=17cm, trim=173 100 137 80, clip=true]{graficas/azotea/disponibilidad}\\
%     	No se encontró perdida en la información en el período descargado [\Diaini/\Mesini/\Anoini~al \Diafin/\Mesfin/\Anofin], sin embargo, se detectaron algunas interrupciones internas dentro del período descargado.\\
% 	\end{tblr}
% \end{table}

\subsection{ESTADO FÍSICO DEL ACELERÓGRAFO}
\begin{table}[H]
    \begin{tblr}{
        row{2} = {bg=SuCielo},
        colspec = {XX}
    }
        \SetCell[c=2]{l} ESTADO FÍSICO DEL ACELERÓGRAFO \\
        \SetCell[c=2]{l} ACELERÓGRAFO EN TERRENO NATURAL \\
        \SetCell{c} \includegraphics[angle=-90, width=7.5cm]{example-image} &
        \SetCell{c} \includegraphics[angle=-90, width=7.5cm]{example-image} \\
        \SetCell{c} \includegraphics[angle=-90, width=7.5cm]{example-image} &
        \SetCell{c} \includegraphics[angle=-90, width=7.5cm]{example-image} \\
        \SetCell[c=2]{l, 15cm} El acelerógrafo se encontró en buenas condiciones, con todos sus perífericos operando con normalidad: sistema de energía, receptor GPS, pantalla LCD de notificación de estado y memoria USB extraíble.
    \end{tblr}
\end{table}

% \begin{table}[H]
%     \begin{tblr}{
%         row{2,Y} = {bg=SuCielo},
%         row{5} = {bg=gray!25},
%         colspec = {XX}
%     }
%         \SetCell[c=2]{l} ESTADO FÍSICO DEL ACELERÓGRAFO\\
%         \SetCell[c=2]{l} ACELERÓGRAFO AZOTEA\\
%         \SetCell{c} \includegraphics[angle=-90, width=7.5cm]{fotos/azotea/acelerografo/1} &
%         \SetCell{c} \includegraphics[angle=-90, width=7.5cm]{fotos/azotea/acelerografo/2} \\
%         \SetCell{c} \includegraphics[angle=-90, width=7.5cm]{fotos/azotea/acelerografo/3} &
%         \SetCell{c} \includegraphics[angle=-90, width=7.5cm]{fotos/azotea/acelerografo/4} \\
%         \SetCell[c=2]{l, 15cm} El acelerógrafo se encontró en buenas condiciones, con todos sus perífericos operando con normalidad: sistema de energía, receptor GPS, pantalla LCD de notificación de estado y memoria USB extraíble.
%     \end{tblr}
% \end{table}


\begin{table}[H]
    \begin{tblr}{
        row{2} = {bg=SuCielo},
        colspec = {XX}
    }
        \SetCell[c=2]{l} COMPONENTES DE LA ESTACIÓN \\
        \SetCell[c=2]{l} ESTACIÓN ACELEROMÉTRICA EN TERRENO NATURAL \\
        \SetCell{c} \includegraphics[angle=-90, width=7.5cm]{example-image} &
        \SetCell{c} \includegraphics[angle=-90, width=7.5cm]{example-image} \\
        \SetCell{c} \includegraphics[angle=-90, width=7.5cm]{example-image} &
        \SetCell{c} \includegraphics[angle=-90, width=7.5cm]{example-image} \\
        \SetCell[c=2]{l, 12cm} Los componentes del sistema o estación sísmica se encuentran en buen estado físico, es decir, cajas de paso, tablero de distribución y componentes internos.
    \end{tblr}
\end{table}

% \begin{table}[H]
%     \begin{tblr}{
%         row{2} = {bg=RoyalBlue!50},
%         row{5} = {bg=gray!25},
%         colspec = {XX}
%     }
%         \SetCell[c=2]{l} COMPONENTES DE LA ESTACIÓN\\
%         \SetCell[c=2]{l} ESTACIÓN ACELEROMÉTRICA AZOTEA\\
%         \SetCell{c} \includegraphics[angle=-90, width=7.5cm]{fotos/azotea/estacion/1} &
%         \SetCell{c} \includegraphics[angle=-90, width=7.5cm]{fotos/azotea/estacion/2} \\
%         \SetCell{c} \includegraphics[angle=-90, width=7.5cm]{fotos/azotea/estacion/3} &
%         \SetCell{c} \includegraphics[angle=-90, width=7.5cm]{fotos/azotea/estacion/4} \\
%         \SetCell[c=2]{l, 12cm} Los componentes del sistema o estación sísmica se encuentran en buen estado físico, es decir, cajas de paso, tablero de distribución y componentes internos.
%     \end{tblr}
% \end{table}

\subsection{ESTADO DEL SISTEMA DE SUMINISTRO ELÉCTRICO}
\begin{longtblr}{
    row{2} = {bg=SuCielo},
    row{3} = {bg=gray!25},
    colspec = {XX}
}
    \SetCell[c=2]{l} SISTEMA DE SUMINISTRO ELÉCTRICO \\
    \SetCell[c=2]{l} TABLERO EN TERRENO NATURAL \\
    \SetCell[c=2]{l} COMPONENTES DEL SISTEMA Y MEDICIONES\\
    \SetCell{c} \includegraphics[angle=-90, width=7.5cm]{example-image} &
    \SetCell{c} \includegraphics[angle=-90, width=7.5cm]{example-image} \\
    \SetCell{c} \includegraphics[angle=-90, width=7.5cm]{example-image} &
    \SetCell{c} \includegraphics[angle=-90, width=7.5cm]{example-image} \\
    \SetCell{c} \includegraphics[angle=-90, width=7.5cm]{example-image} &
    \SetCell{c} \includegraphics[angle=-90, width=7.5cm]{example-image} \\
    \SetCell[c=2]{l} Medición de parámetros eléctricos en la entrada del rectificador y controlador de carga.
\end{longtblr}

% \begin{table}[H]
%     \begin{tblr}{
%         row{2} = {bg=RoyalBlue!50},
%         row{3,6} = {bg=gray!25},
%         colspec = {XX}
%     }
%         \SetCell[c=2]{l} SISTEMA DE SUMINISTRO ELÉCTRICO\\
%         \SetCell[c=2]{l} TABLERO AZOTEA\\
%         \SetCell[c=2]{l} 1. COMPONENTES DEL SISTEMA Y MEDICIONES\\
%         \SetCell{c} \includegraphics[angle=-90, width=7.5cm]{fotos/azotea/tablero/1} &
%         \SetCell{c} \includegraphics[angle=-90, width=7.5cm]{fotos/azotea/tablero/2} \\
%         \SetCell{c} \includegraphics[angle=-90, width=7.5cm]{fotos/azotea/tablero/3} &
%         \SetCell{c} \includegraphics[angle=-90, width=7.5cm]{fotos/azotea/tablero/4} \\
%         \SetCell[c=2]{l} Medición de parámetros eléctricos en la entrada del rectificador y controlador de carga.
%     \end{tblr}
% \end{table}
% \begin{table}[H]
%     \begin{tblr}{
%         row{2} = {bg=RoyalBlue!50},
%         row{3,5} = {bg=gray!25},
%         colspec = {XX}
%     }
%         \SetCell[c=2]{l} SISTEMA DE SUMINISTRO ELÉCTRICO\\
%         \SetCell[c=2]{l} TABLERO AZOTEA\\
%         \SetCell[c=2]{l} 2. COMPONENTES DEL SISTEMA Y MEDICIONES\\
%         \SetCell{c} \includegraphics[angle=-90, width=7.5cm]{fotos/azotea/tablero/5} &
%         \SetCell{c} \includegraphics[angle=-90, width=7.5cm]{fotos/azotea/tablero/6} \\
%         \SetCell[c=2]{l} Prueba de descarga de batería.
%     \end{tblr}
% \end{table}


\subsection{ESTADO DEL SENSOR INTERNO}
\begin{table}[H]
    \begin{tblr}{
        row{2} = {bg=SuCielo},
        row{3} = {bg=gray!25},
        colspec = {X}
    }
        ESTADO DEL SENSOR INTERNO \\
        ACELERÓGRAFO EN TERRENO NATURAL \\
        1. CENTRADO DE MASA DE CADA CANAL - RETIRO DE OFFSET \\
        % \SetCell{c}\includegraphics[angle=-90, width=17cm, trim=2200 900 1310 1010, clip=true]{example-image} \\
        % \SetCell{c}\includegraphics[angle=-90, width=17cm, trim=800 0 1175 0, clip=true]{example-image} \\
        {Superior: Offset Inicial (mg) \\ Inferior: Offset Final (ug)}
    \end{tblr}
\end{table}
\begin{table}[H]
    \begin{tblr}{
        row{2} = {bg=SuCielo},
        row{3} = {bg=gray!25},
        colspec = {X}
    }
        ESTADO DEL SENSOR INTERNO \\
        ACELERÓGRAFO EN TERRENO NATURAL \\
        2. EXCITACIÓN CON GOLPES CERCA DE LA BASE DE CONCRETO \\
        % \SetCell{c} \includegraphics[angle=0, width=17cm, trim=120 160 900 400, clip=true, draft=false]{graficas/sotano/offset/1} \\
        % \SetCell{c} \includegraphics[angle=0, width=17cm, trim=120 160 900 400, clip=true, draft=false]{graficas/sotano/offset/2} \\
        {Superior: Formas de onda del ruido de fondo \\ Inferior: Excitación de los sensores con golpes sobre la base de concreto}
    \end{tblr}
\end{table}
\begin{table}[H]
    \begin{tblr}{
        row{2} = {bg=SuCielo},
        row{3} = {bg=gray!25},
        colspec = {X}
    }
        ESTADO DEL SENSOR INTERNO \\
        ACELERÓGRAFO EN TERRENO NATURAL \\
        3. CALIBRACIÓN DEL ACELERÓGRAFO - INYECCIÓN DE LA SEÑAL DE CALIBRACIÓN [OUTPUT - Count] \\
        % \SetCell{c} \includegraphics[angle=0, width=17cm, trim=120 160 900 400, clip=true, draft=false]{graficas/sotano/estimulo} \\
        {Superior: Señal de calibración inyectada vía software desde el acelerógrafo (forma de onda: sinuosidal de Amplitud: 1VDC y Frecuencia: 0.25Hz) \\ Inferior: Señal censada y registrada por el acelerógrafo. Se valida la calibración reconstruyendo la señal registrada en términos de cuentas.}
    \end{tblr}
\end{table}
\begin{table}[H]
    \begin{tblr}{
        row{2} = {bg=SuCielo},
        row{3} = {bg=gray!25},
        colspec = {X}
    }
        ESTADO DEL SENSOR INTERNO \\
        ACELERÓGRAFO EN TERRENO NATURAL \\
        4. CALIBRACIÓN DEL ACELERÓGRAFO - INYECCIÓN DE LA SEÑAL DE CALIBRACIÓN [OUTPUT - \unit{\centi\metre\per\second\squared}] \\
        % \SetCell{c} \includegraphics[angle=0, width=17cm, trim=120 160 900 400, clip=true, draft=false]{graficas/sotano/calibracion}\\
        % \SetCell{c} \includegraphics[angle=0, width=17cm, trim=70 30 20 60, clip=true, draft=false]{graficas/sotano/descompuesta/cuentas}\\
        \SetCell{l, 16cm} {Superior: Señal de calibración inyectada vía software desde el acelerógrafo (forma de onda: pulso cuadrado y 5 seg de duración)\\Inferior: Señal sensada y registrada por el acelerógrafo (en cuentas)}\\
    \end{tblr}
\end{table}
\begin{table}[H]
    \begin{tblr}{
        row{2} = {bg=SuCielo},
        row{3} = {bg=gray!25},
        colspec = {X}
    }
        ESTADO DEL SENSOR INTERNO \\
        ACELERÓGRAFO EN TERRENO NATURAL \\
        5. CALIBRACIÓN DEL ACELERÓGRAFO - INYECCIÓN DE LA SEÑAL DE CALIBRACIÓN [ESPECTROGRAMA] \\
        % \SetCell{c} \includegraphics[angle=0, width=17cm, trim=90 35 85 60, clip=true, draft=false]{graficas/sotano/espectrograma}\\
        {
            Espectrograma de la señal calibrada \\
            Eje X: Señal de calibración en el dominio del tiempo, se observa que en la frecuencia de calibración los valores de amplitud en términos de cuentas. \\
            Eje Y: Señal de calibración en el dominio de la frecuencia, se observa que la frecuencia de calibración es de 0.25Hz, validando la reconstrucción de la señal.
        } \\
    \end{tblr}
\end{table}


% \begin{table}[H]
%     \begin{tblr}{
%         row{2} = {bg=RoyalBlue!50},
%         row{3,5} = {bg=gray!25},
%         colspec = {X}
%     }
%         ESTADO DEL SENSOR INTERNO\\
%         ACELERÓGRAFO SÓTANO\\
%         6. CALIBRACIÓN DEL ACELERÓGRAFO - INYECCIÓN DE LA SEÑAL DE CALIBRACIÓN [ESPECTROGRAMA]\\
%         \SetCell{c} \includegraphics[angle=0, width=17cm, trim=90 35 85 60, clip=true, draft=false]{graficas/sotano/espectrograma}\\
%         {Espectrograma de la señal calibrada\\Lateral izquierda: Señal de calibración en el dominio de la frecuencia\\Inferior derecha: Señal de calibración en el dominio del tiempo}\\
%     \end{tblr}
% \end{table}

% \begin{table}[H]
%     \begin{tblr}{
%         row{2} = {bg=RoyalBlue!50},
%         row{3,6} = {bg=gray!25},
%         colspec = {X}
%     }
%         ESTADO DEL SENSOR INTERNO\\
%         ACELERÓGRAFO AZOTEA\\
%         1. NIVELACIÓN MECÁNICA EN LA HORIZONTAL\\
%         \SetCell{c} \includegraphics[angle=-90, width=17cm, trim=2200 900 1310 1010, clip=true]{fotos/azotea/nivelacion/1} \\
%         \SetCell{c} \includegraphics[angle=-90, width=17cm, trim=800 0 1175 0, clip=true]{fotos/azotea/nivelacion/2} \\
%         Se niveló la horizontal del acelerógrafo mediante el perno de anclaje y burbuja de nivel.
%     \end{tblr}
% \end{table}
% \begin{table}[H]
%     \begin{tblr}{
%         row{2} = {bg=RoyalBlue!50},
%         row{3,6} = {bg=gray!25},
%         colspec = {X}
%     }
%         ESTADO DEL SENSOR INTERNO\\
%         ACELERÓGRAFO AZOTEA\\
%         2. NIVEL DE OFFSET DE CADA CANAL ANTES Y DESPUES DEL MANTENIMIENTO \\
%         \SetCell{c} \includegraphics[angle=0, width=17cm, trim=120 160 900 400, clip=true, draft=false]{graficas/azotea/offset/1} \\
%         \SetCell{c} \includegraphics[angle=0, width=17cm, trim=120 160 900 400, clip=true, draft=false]{graficas/azotea/offset/2} \\
%         {Superior: Offset antes de realizar el mantenimiento\\Inferior: Offset luego de realizar el mantenimiento}
%     \end{tblr}
% \end{table}
% \begin{table}[H]
%     \begin{tblr}{
%         row{2} = {bg=RoyalBlue!50},
%         row{3,5} = {bg=gray!25},
%         colspec = {X}
%     }
%         ESTADO DEL SENSOR INTERNO\\
%         ACELERÓGRAFO AZOTEA\\
%         3. ESTIMULACIÓN CON GOLPES CERCA DE LA BASE DE CONCRETO\\
%         \SetCell{c} \includegraphics[angle=0, width=17cm, trim=120 160 900 400, clip=true, draft=false]{graficas/azotea/estimulo} \\
%         Se estimuló el sensor con pequeños golpes a los costados para validar su funcionamiento en los tres canales.
%     \end{tblr}
% \end{table}
% \begin{table}[H]
%     \begin{tblr}{
%         row{2} = {bg=RoyalBlue!50},
%         row{3,6} = {bg=gray!25},
%         colspec = {X}
%     }
%         ESTADO DEL SENSOR INTERNO\\
%         ACELERÓGRAFO AZOTEA\\
%         4. CALIBRACIÓN DEL ACELERÓGRAFO - INYECCIÓN DE LA SEÑAL DE CALIBRACIÓN [OUTPUT - count]\\
%         \SetCell{c} \includegraphics[angle=0, width=17cm, trim=120 160 900 400, clip=true, draft=false]{graficas/azotea/calibracion}\\
%         \SetCell{c} \includegraphics[angle=0, width=17cm, trim=70 30 20 60, clip=true, draft=false]{graficas/azotea/descompuesta/cuentas}\\
%         \SetCell{l, 16cm} {Superior: Señal de calibración inyectada vía software desde el acelerógrafo (forma de onda: pulso cuadrado y 5 seg de duración)\\Inferior: Señal sensada y registrada por el acelerógrafo (en cuentas)}\\
%     \end{tblr}
% \end{table}
% \begin{table}[H]
%     \begin{tblr}{
%         row{2} = {bg=RoyalBlue!50},
%         row{3,6} = {bg=gray!25},
%         colspec = {X}
%     }
%         ESTADO DEL SENSOR INTERNO\\
%         ACELERÓGRAFO AZOTEA\\
%         5. CALIBRACIÓN DEL ACELERÓGRAFO - INYECCIÓN DE LA SEÑAL DE CALIBRACIÓN [OUTPUT - mm/s2]\\
%         \SetCell{c} \includegraphics[angle=0, width=17cm, trim=120 160 900 400, clip=true, draft=false]{graficas/azotea/calibracion}\\
%         \SetCell{c} \includegraphics[angle=0, width=17cm, trim=70 30 20 60, clip=true, draft=false]{graficas/azotea/descompuesta/aceleracion}\\
%         {Superior: Señal de calibración inyectada vía software desde el acelerógrafo (forma de onda: pulso cuadrado y 5 seg de duración)\\Inferior: Señal reconstruída utilizando la respuesta instrumental (en términos de aceleración)}\\
%     \end{tblr}
% \end{table}
% \begin{table}[H]
%     \begin{tblr}{
%         row{2} = {bg=RoyalBlue!50},
%         row{3,5} = {bg=gray!25},
%         colspec = {X}
%     }
%         ESTADO DEL SENSOR INTERNO\\
%         ACELERÓGRAFO AZOTEA\\
%         6. CALIBRACIÓN DEL ACELERÓGRAFO - INYECCIÓN DE LA SEÑAL DE CALIBRACIÓN [ESPECTROGRAMA]\\
%         \SetCell{c} \includegraphics[angle=0, width=17cm, trim=90 35 85 60, clip=true, draft=false]{graficas/azotea/espectrograma}\\
%         {Espectrograma de la señal calibrada\\Lateral izquierda: Señal de calibración en el dominio de la frecuencia\\Inferior derecha: Señal de calibración en el dominio del tiempo}\\
%     \end{tblr}
% \end{table}


% \subsection{ESTADO FINAL DEL ACELERÓGRAFO}
% \begin{table}[H]
%     \begin{tblr}{
%         row{2} = {bg=RoyalBlue!50},
%         row{4} = {bg=gray!25},
%         colspec = {X}
%     }
%         STATUS GLOBAL DEL ACELERÓGRAFO\\
%         ESTACIÓN SÓTANO\\
%         \SetCell{c} \includegraphics[angle=0, width=17cm, trim=120 300 1200 190, clip=true, draft=false]{graficas/sotano/status} \\
%         Status global de los acelerógrafos en sótano al culminar el mantenimiento - Estado final: Estación operativa
%     \end{tblr}
% \end{table}

% \begin{table}[H]
%     \begin{tblr}{
%         row{2} = {bg=RoyalBlue!50},
%         row{4} = {bg=gray!25},
%         colspec = {X}
%     }
%         STATUS GLOBAL DEL ACELERÓGRAFO\\
%         ESTACIÓN AZOTEA\\
%         \SetCell{c} \includegraphics[angle=0, width=17cm, trim=120 300 1200 190, clip=true, draft=false]{graficas/azotea/status} \\
%         Status global de los acelerógrafos en azotea al culminar el mantenimiento - Estado final: Estación operativa
%     \end{tblr}
% \end{table}